\documentclass[a4paper,11pt]{article}

% =============================================================================
%  Charte graphique des documents LaTeX du projet (gabarit commenté).
%  Le préambule fixe la mise en forme ; le corps montre chaque élément une fois
%  et se remplace par le contenu du document. Les passages propres au projet
%  sont marqués « À ADAPTER ».
%  Version sans : babel-french, caption, csquotes, dsfont, enumitem, float,
%  mathrsfs, mdframed, newtxtext, pdflscape, tikz/pgf, titlesec, tocloft.
%  Moteur : pdfLaTeX. Compatible avec les distributions anciennes
%  (MiKTeX 2.9, TeX Live 2012 et suivantes) : aucun crochet du noyau récent.
% =============================================================================

\usepackage[utf8]{inputenc}
\usepackage[T1]{fontenc}
% \headheight est laissé à 12 pt : fancyhdr le signale et l'agrandit lui-même
% (avertissement « \headheight is too small »), exactement comme dans le
% document source ; le fixer ici décalerait l'en-tête ou le texte de 1,6 pt.
\usepackage[left=2.5cm, right=2.5cm, top=2.5cm, bottom=2.5cm]{geometry}
\usepackage{setspace}
\usepackage{graphicx}
\usepackage{etoolbox}
\usepackage{keyval}
\usepackage{fancyhdr}
\usepackage{amsmath,amssymb,amsfonts}
\usepackage{url}
\usepackage[table]{xcolor}  % [table] : \rowcolors sur les anciennes versions
\usepackage{array}
\usepackage{booktabs}
\usepackage{colortbl}
\usepackage{longtable}
\usepackage{lscape}   % remplace pdflscape (rotation à l'écran : voir plus bas)
\usepackage{framed}   % remplace mdframed (encadré sécable)
\usepackage{atbegshi}  % crochets de sortie de page (remplace \AddToHook)
\usepackage[unicode]{hyperref}

% ------------------------------------------------------------- polices -------
% newtxtext composait le texte en TeX Gyre Termes X et le code en txtt ; les
% mathématiques restaient en Computer Modern. On retrouve la même chose sans
% newtxtext : TeX Gyre Termes (tgtermes) pour le texte, repli sur le Times de
% PSNFSS (ptm) s'il est absent ; txtt (txfonts) pour \ttfamily.
\IfFileExists{tgtermes.sty}
  {\usepackage{tgtermes}}
  {\renewcommand{\rmdefault}{ptm}}
\IfFileExists{t1txtt.fd}{\renewcommand{\ttdefault}{txtt}}{}

% ---------------------------------------------------------------- charte ----
% Bleu financier pour la structure, vert pour l'accent et les sorties, deux
% gris de tableau.
\definecolor{maincolor}{RGB}{0,70,140}
\definecolor{greenline}{RGB}{0,180,80}
\definecolor{tableheader}{RGB}{200,220,250}
\definecolor{tablerow}{RGB}{245,245,245}
\definecolor{financeblue}{RGB}{0,70,140}

% À ADAPTER : métadonnées du document (page de garde, pied de page, PDF).
\newcommand{\annee}{2026}
\newcommand{\datedoc}{\today}
\newcommand{\versiondoc}{1.0}
\newcommand{\statutdoc}{Projet}
\newcommand{\titredoc}{Titre court du document}

\hypersetup{colorlinks=true, linkcolor=financeblue, urlcolor=financeblue,
            citecolor=financeblue, pdftitle={Titre du document}}

\makeatletter

% ================================================ français sans babel ========
% Reprise, à la main, des effets de babel-french visibles dans le document
% source.
% (1) Césure française (motifs chargés dans le format pdfLaTeX).
\ifdefined\l@french
  \language=\l@french
\else
  \PackageWarningNoLine{charte}{Motifs de césure français absents :
    césure anglaise utilisée}
\fi
\lefthyphenmin=2 \righthyphenmin=2
% (2) Espacement « à la française » après les points.
\frenchspacing
% (3) Libellés.
\renewcommand{\contentsname}{Table des matières}
\renewcommand{\listfigurename}{Table des figures}
\renewcommand{\listtablename}{Liste des tableaux}
\renewcommand{\refname}{Références}
\renewcommand{\indexname}{Index}
\renewcommand{\figurename}{Figure}
\renewcommand{\tablename}{Table}
\renewcommand{\partname}{Partie}
\renewcommand{\appendixname}{Annexe}
\renewcommand{\abstractname}{Résumé}
% (4) Date en français (« 1er janvier 2026 »).
\renewcommand{\today}{\number\day\ifnum\day=1 \textsuperscript{er}\fi\space
  \ifcase\month\or janvier\or février\or mars\or avril\or mai\or juin\or
  juillet\or août\or septembre\or octobre\or novembre\or décembre\fi
  \space\number\year}
% (5) Premier paragraphe indenté après un titre.
\let\@afterindentfalse\@afterindenttrue
\@afterindenttrue
% (6) Listes : espacement vertical resserré de babel-french, appliqué à tous
%     les environnements construits sur \list.
\let\charte@listORI\list
\newdimen\charte@parskip
\def\charte@listV{%
  \setlength{\topsep}{0.8ex plus 0.4ex minus 0.4ex}%
  \setlength{\partopsep}{0.4ex plus 0.2ex minus 0.2ex}%
  \charte@parskip=\parskip
  \addtolength{\topsep}{-\charte@parskip}%
  \addtolength{\partopsep}{\charte@parskip}%
  \setlength{\itemsep}{0.4ex plus 0.2ex minus 0.2ex}%
  \setlength{\parsep}{0.4ex plus 0.2ex minus 0.2ex}%
  \ifdim\charte@parskip>\z@
    \setlength{\parsep}{\charte@parskip}%
    \addtolength{\itemsep}{-\charte@parskip}%
  \fi}
\def\list#1#2{\charte@listORI{#1}{\charte@listV #2}}
% (7) Ponctuation haute : sans babel-french, l'espace qui précède « : » et
%     l'espace fine qui précède « ; ! ? » ne sont plus ajoutées. On saisit
%     « ~: » (espace insécable) et « \fsp; », « \fsp! », « \fsp? » (espace
%     fine insécable, de même largeur que celle de babel-french).
\newcommand{\fsp}{\unskip\nobreak\hskip .5\fontdimen2\font\relax}
% (8) Puces en tiret cadratin à tous les niveaux.
\renewcommand{\labelitemi}{\textemdash}
\renewcommand{\labelitemii}{\textemdash}
\renewcommand{\labelitemiii}{\textemdash}
\renewcommand{\labelitemiv}{\textemdash}

% ================================================ listes sans enumitem =======
% itemize et enumerate acceptent un argument optionnel à clés, dont la
% syntaxe est celle d'enumitem pour les clés usuelles :
%   \begin{itemize}[itemsep=3pt] ... \end{itemize}
% Clés reconnues : itemsep, topsep, parsep, partopsep, leftmargin,
% labelsep, listparindent.
\define@key{charteliste}{itemsep}{\setlength{\itemsep}{#1}}
\define@key{charteliste}{topsep}{\setlength{\topsep}{#1}}
\define@key{charteliste}{parsep}{\setlength{\parsep}{#1}}
\define@key{charteliste}{partopsep}{\setlength{\partopsep}{#1}}
\define@key{charteliste}{leftmargin}{\setlength{\leftmargin}{#1}}
\define@key{charteliste}{labelsep}{\setlength{\labelsep}{#1}}
\define@key{charteliste}{listparindent}{\setlength{\listparindent}{#1}}

\renewenvironment{itemize}[1][]{%
  \ifnum\@itemdepth>\thr@@\@toodeep\else
    \advance\@itemdepth\@ne
    \edef\@itemitem{labelitem\romannumeral\the\@itemdepth}%
    \expandafter\list\csname\@itemitem\endcsname
      {\def\makelabel##1{\hss\llap{##1}}\setkeys{charteliste}{#1}}%
  \fi}{\endlist}

\renewenvironment{enumerate}[1][]{%
  \ifnum\@enumdepth>\thr@@\@toodeep\else
    \advance\@enumdepth\@ne
    \edef\@enumctr{enum\romannumeral\the\@enumdepth}%
    \expandafter\list\csname label\@enumctr\endcsname
      {\usecounter\@enumctr\def\makelabel##1{\hss\llap{##1}}%
       \setkeys{charteliste}{#1}}%
  \fi}{\endlist}

% ================================================ légendes sans caption ======
% Étiquette en gras, couleur maincolor, petites capitales ; séparateur
% « : » précédé d'une espace insécable, comme dans le document source.
\renewcommand{\fnum@figure}{{\scshape\figurename}\nobreakspace\thefigure}
\renewcommand{\fnum@table}{{\scshape\tablename}\nobreakspace\thetable}
\newcommand{\charte@etiquette}[1]{%
  {\bfseries\color{maincolor}#1\nobreakspace:}\space}
% Comme avec caption : légende composée en interligne simple, entre deux
% étais, sans ressort interligne au-dessus ; centrée si elle tient sur une
% ligne, justifiée sinon.
\long\def\@makecaption#1#2{%
  \vskip\abovecaptionskip
  \sbox\@tempboxa{\normalfont\charte@etiquette{#1}#2}%
  \nointerlineskip
  \ifdim\wd\@tempboxa>\hsize
    \hb@xt@\hsize{\vbox{\setstretch{1}\normalfont
      \parindent\z@ \leftskip\z@ \rightskip\z@ \parfillskip\z@\@plus1fil
      \strut\charte@etiquette{#1}#2\strut\par}}%
  \else
    \global\@minipagefalse
    \hb@xt@\hsize{\hfil\setstretch{1}\normalfont\strut\box\@tempboxa\hfil}%
  \fi
  \vskip\belowcaptionskip}

% Figures et tableaux « ici » : reprise exacte de la position [H] du
% package float (boîte non flottante, \intextsep avant et après).
\newsavebox\charte@boiteH
\def\charte@debutH#1{%
  \def\@captype{#1}%
  \global\setbox\charte@boiteH\color@vbox\normalcolor
    \vbox\bgroup\hsize\columnwidth\@parboxrestore
      \@floatboxreset\@setnobreak\ignorespaces}
\def\charte@finH{%
  \@endfloatbox
  \vskip\intextsep\box\charte@boiteH\vskip\intextsep\relax}
\newenvironment{figureH}{\charte@debutH{figure}}{\charte@finH}
\newenvironment{tableH}{\charte@debutH{table}}{\charte@finH}

% ================================================ en-têtes et pieds ==========
\renewcommand{\baselinestretch}{1.08}
\pagestyle{fancy}
\fancyhf{}
\fancyhead[L]{\color{maincolor}\nouppercase{\leftmark}}
\renewcommand{\headrulewidth}{0pt}
\renewcommand{\footrulewidth}{0.4pt}
\fancyfoot[R]{\color{greenline}\textbf{page \thepage}}
\fancyfoot[L]{\color{maincolor}\titredoc}
\fancypagestyle{plain}{\fancyhf{}%
  \renewcommand{\headrulewidth}{0pt}\renewcommand{\footrulewidth}{0.4pt}%
  \fancyfoot[R]{\color{greenline}\textbf{page \thepage}}%
  \fancyfoot[L]{\color{maincolor}\titredoc}}

% ================================================ titres sans titlesec =======
% Bandeau : filet bleu pleine largeur, titre vert en dessous. Utilisé seul
% pour les pages non numérotées (objet, sommaire) et par \section.
\newcommand{\titrebandeau}[1]{%
  \noindent{\color{financeblue}\rule{\linewidth}{2.5pt}}
  \par\vspace{5mm}
  {\color{greenline}\raggedright\fontsize{19}{23}\bfseries\selectfont #1\par}
  \vspace{6mm}}

% Numéros de titre : aucun pour \section (le numéro ne figure qu'au
% sommaire et dans l'en-tête), 0,7 em puis 0,6 em d'écart ensuite.
\def\@seccntformat#1{\csname charte@num@#1\endcsname}
\def\charte@num@section{}
\def\charte@num@subsection{\strut\thesubsection\hspace{0.7em}}
\def\charte@num@subsubsection{\strut\thesubsubsection\hspace{0.6em}}
\def\charte@num@paragraph{\theparagraph\quad}
\def\charte@num@subparagraph{\thesubparagraph\quad}

% \@startsection passe le titre, entre accolades, à la dernière commande de
% son argument de style. Pour \section, cette commande est \charte@bandeau :
% une ligne vide (étai) au-dessus du bandeau et une autre en dessous, et un
% étai dans chaque ligne, comme la forme « display » de titlesec avec
% étiquette vide.
\newcommand{\charte@bandeau}[1]{%
  \parindent\z@ \everypar{\strut}%
  {\strut\par}\nobreak
  \titrebandeau{#1}\nobreak
  \kern\z@\strut\par}
% Titres en ligne : un étai en fin de titre (le premier est dans le numéro).
\newcommand{\charte@etai}[1]{\charte@etai@#1}
\def\charte@etai@#1\@@par{#1\kern\z@\strut\@@par}
\renewcommand{\section}{\@startsection{section}{1}{\z@}%
  {\z@}{1sp}{\normalfont\charte@bandeau}}
% Espacements : ceux de titlesec, soit 3,25 ex (+1 ex, -0,2 ex) avant et
% 1,5 ex (+0,2 ex) après, l'ex étant celui que titlesec mémorise au
% chargement (Computer Modern à 11 pt, soit 4,71341 pt).
\newdimen\charte@ex \charte@ex=4.71341pt
\newskip\charte@avanttitre
\charte@avanttitre=-3.25\charte@ex \@plus -1\charte@ex \@minus -.2\charte@ex
\newskip\charte@aprestitre
\charte@aprestitre=1.5\charte@ex \@plus .2\charte@ex
\renewcommand{\subsection}{\@startsection{subsection}{2}{\z@}%
  {\charte@avanttitre}{\charte@aprestitre}%
  {\normalfont\large\bfseries\color{financeblue}\charte@etai}}
\renewcommand{\subsubsection}{\@startsection{subsubsection}{3}{\z@}%
  {\charte@avanttitre}{\charte@aprestitre}%
  {\normalfont\normalsize\bfseries\color{financeblue}\charte@etai}}

% ================================================ sommaire sans tocloft ======
% Sections : gras bleu, numéro sur 2,2 em, numéro de page bleu (maigre),
% sans points de conduite. Sous-sections : retrait 1,5 em, numéro sur 2,3 em,
% points de conduite. L'em est celui que tocloft mémorisait au chargement
% (10,88788 pt), pour un sommaire identique au document source.
\newdimen\charte@em \charte@em=10.88788pt
\renewcommand*{\l@section}[2]{%
  \ifnum\c@tocdepth>\z@
    \addpenalty\@secpenalty
    \addvspace{1.0\charte@em \@plus\p@}%
    \setlength\@tempdima{2.2\charte@em}%
    \begingroup
      \parindent\z@ \rightskip\@pnumwidth
      \parfillskip -\@pnumwidth
      \leavevmode
      \advance\leftskip\@tempdima
      \hskip -\leftskip
      {\bfseries\color{financeblue}#1}\nobreak\hfil\nobreak
      \hb@xt@\@pnumwidth{\hss\color{financeblue}#2}\par
    \endgroup
  \fi}
\renewcommand*{\l@subsection}{\@dottedtocline{2}{1.5\charte@em}{2.3\charte@em}}
\setcounter{tocdepth}{2}

% ================================================ environnements de texte ====
\newcommand{\code}[1]{\texttt{\small #1}}
\newcommand{\refl}[1]{{\small\textcolor{maincolor!70!black}{\textit{Référence.} #1}}\par\smallskip}

% Encadré de citation réglementaire : filet vertical bleu en marge.
\newcommand{\reglement}[2]{%
  \par\vspace{2mm}\noindent\begin{minipage}{\linewidth}
  \small\setlength{\parindent}{0pt}
  \textcolor{financeblue}{\vrule width 2pt}\hspace{4mm}%
  \begin{minipage}{\dimexpr\linewidth-7mm\relax}
  \textcolor{financeblue}{\textbf{#1}}
  \par\vspace{2pt}
  {\itshape #2}
  \end{minipage}
  \end{minipage}\par\vspace{3mm}}

% Encadré générique : fond gris, marge intérieure 6 mm à gauche et 10 pt à
% droite, 4 pt en haut et en bas, 3 mm avant (sauf après un titre) et
% après ; premier paragraphe non
% indenté. Sécable entre pages (framed).
\newenvironment{encadre}{%
  \def\FrameCommand##1{{\setlength{\fboxsep}{\z@}%
    \colorbox{tablerow}{\hspace*{6mm}%
      \vbox{\vskip4pt\setbox\@tempboxa\hbox{##1}\dp\@tempboxa\z@
            \box\@tempboxa\vskip4pt}\hspace*{10pt}}}}%
  \@ifundefined{OuterFrameSep}{}{\setlength{\OuterFrameSep}{\z@}}%
  \topsep\z@ \partopsep\z@ % framed antérieur à 0.96 : espacement via \topsep
  \par
  \if@nobreak\else\vskip 3mm\fi
  \nointerlineskip
  \MakeFramed{\advance\hsize-\width \FrameRestore
    \everypar{{\setbox\z@\lastbox}\everypar{}}}}%
  {\endMakeFramed\par
   \vskip 3mm\prevdepth 4pt
   \global\@nobreakfalse\global\@endpetrue}

% Fiches : six rubriques fixes. À ADAPTER : intitulés des rubriques (ici,
% ceux d'une fiche de test statistique).
\newcounter{ficheid}
\newenvironment{fiche}[1]{%
  \stepcounter{ficheid}%
  \phantomsection
  \subsection{#1}%
  \label{fiche:\theficheid}%
}{%
  \end{rubriques}\vspace{0.5em}%
}
% Liste des rubriques : aucun retrait ; l'intitulé fait partie du
% paragraphe (ses espaces se justifient avec la ligne, comme le style
% « unboxed » d'enumitem), suivi d'un écart fixe \labelsep.
\newenvironment{rubriques}{%
  \list{}{\setlength{\leftmargin}{\z@}\setlength{\labelwidth}{\z@}%
          \setlength{\itemindent}{\z@}%
          \setlength{\itemsep}{3pt}\setlength{\topsep}{3pt}%
          \setlength{\parsep}{2pt}%
          \let\makelabel\@gobble}}{\endlist}
\newcommand{\rubrique}[1]{%
  \item[]\textcolor{financeblue}{\textbf{#1}}\hskip\labelsep\ignorespaces}
\newcommand{\Cadre}{%
  \begin{rubriques}%
  \rubrique{1. Cadre, notations et hypothèses de validité (non testées).}}
\newcommand{\Hzero}{\rubrique{2. Hypothèse nulle.}}
\newcommand{\Stat}{\rubrique{3. Statistique de test.}}
\newcommand{\Loi}{\rubrique{4. Loi sous $H_0$.}}
\newcommand{\Pval}{\rubrique{5. p-value.}}
\newcommand{\Usage}{\rubrique{6. Usage et limites.}}

% Titre de partie : un filet bleu, « Partie » en petites capitales, titre
% vert.
\newcommand{\partie}[2]{%
  \clearpage
  \vspace*{1cm}
  \noindent{\color{financeblue}\rule{\linewidth}{2pt}}\par\vspace{3mm}
  {\color{financeblue}\raggedright\normalsize\bfseries\scshape #1\par}
  \vspace{1.5mm}
  {\color{greenline}\raggedright\fontsize{20}{24}\bfseries\selectfont #2\par}
  \vspace{5mm}
  \addcontentsline{toc}{section}{\protect\numberline{}\textsc{#1} --- #2}}

% ================================================ page de garde sans TikZ ====
% \charte@surpage{<matériel vertical>} superpose, sur la prochaine page
% expédiée, un contenu placé depuis le coin supérieur gauche du papier
% (atbegshi : fonctionne avec les noyaux LaTeX anciens comme récents).
\newcommand{\charte@surpage}[1]{%
  \AtBeginShipoutNext{%
    \global\setbox\AtBeginShipoutBox\vbox{\offinterlineskip
      \vbox to\z@{\kern-1in\kern-\voffset
        \hbox to\z@{\kern-1in\kern-\hoffset
          \vbox to\z@{#1\vss}\hss}\vss}%
      \box\AtBeginShipoutBox}}}
% Trois filets : bleu de 3,2 à 3,28 cm du haut avec amorce verte de 5 cm,
% bleu de 3,6 à 3,53 cm du bas, entre les marges de 2,5 cm.
% \leavevmode ouvre le paragraphe comme le faisait la figure TikZ du document
% source : la mention d'en-tête se place ainsi au-dessus du filet, en retrait.
\newcommand{\filetsgarde}{%
  \charte@surpage{%
    \kern3.2cm
    \moveright2.5cm\hbox{\color{financeblue}%
      \rule{\dimexpr\paperwidth-5cm\relax}{0.08cm}}%
    \kern-0.08cm
    \moveright2.5cm\hbox{\color{greenline}\rule{5cm}{0.08cm}}%
    \kern\dimexpr\paperheight-3.6cm-3.28cm\relax
    \moveright2.5cm\hbox{\color{financeblue}%
      \rule{\dimexpr\paperwidth-5cm\relax}{0.07cm}}}%
  \leavevmode}

% ================================================ pages en paysage ===========
% lscape tourne le contenu ; comme pdflscape, on ajoute /Rotate 90 aux pages
% concernées pour qu'elles s'affichent à l'endroit à l'écran.
\newif\ifcharte@paysage
\apptocmd{\landscape}{\global\charte@paysagetrue}{}{}
\apptocmd{\endlandscape}{\global\charte@paysagefalse}{}{}
\AtBeginShipout{%
  \ifcharte@paysage\global\pdfpageattr{/Rotate 90}%
  \else\global\pdfpageattr{}\fi}

\makeatother

\emergencystretch=1.5em

\begin{document}

% ============================================================ page de garde ==
\thispagestyle{empty}
\filetsgarde
\vspace*{2.1cm}
\noindent
{\color{financeblue}\normalsize\bfseries\scshape Édition \annee}

\vspace{1.7cm}
\noindent
{\color{financeblue}\fontsize{28}{34}\bfseries\selectfont
 Titre principal\\ du document\par}

\vspace{1.1cm}
\noindent
{\color{black}\Large
 Sous-titre sur une première ligne\\
 Sous-titre sur une seconde ligne\par}

\vspace{0.7cm}
\noindent
{\color{black}\normalsize\itshape
 Mention de cadrage --- texte de référence, périmètre ou version du code\par}

\vfill
\noindent
{\footnotesize\color{black}
 Établie le \datedoc\quad\textbar\quad Version \versiondoc\quad\textbar\quad
 Statut~: \statutdoc\quad\textbar\quad Référence \code{chemin/du/fichier}\par}
\vspace{2.6cm}
\clearpage

% ================================================================= objet ====
\titrebandeau{Objet du document}
\markboth{Objet du document}{}

\noindent
Ce gabarit reproduit la \textbf{charte graphique} de la documentation
du projet. Chaque élément de la charte y figure
une fois, dans son contexte d'usage~: page de garde, page d'objet, sommaire,
titres de partie, de section et de sous-section, fiches de test, encadrés,
citations réglementaires, tableaux, tableau long en paysage et figure
légendée.

\medskip\noindent
Il est organisé en trois parties. La \textbf{partie~I} présente la palette,
la typographie et les niveaux de titre. La \textbf{partie~II} montre les
blocs de texte (encadrés, citations, références, listes) et les tableaux. La
\textbf{partie~III} donne le modèle de fiche de test et de figure, puis la
table de correspondance des packages remplacés.

\vspace{8mm}
\reglement{Article 1\textsuperscript{er} du texte de référence}{Extrait
cité du texte de référence, en italique, précédé de son intitulé en gras bleu
et bordé d'un filet vertical en marge.}

\vspace{6mm}
\noindent
\begin{center}
\small
\rowcolors{2}{tablerow}{white}
\begin{tabular}{@{}p{5.2cm}p{9.6cm}@{}}
\rowcolor{tableheader}
\textcolor{financeblue}{\textbf{Élément}} & \textcolor{financeblue}{\textbf{Référence}} \\
Texte de référence & Intitulé du texte, articles concernés\fsp; annexes \\
Cahier des charges & \code{docs/exigences.md} \\
Module de calcul & \code{<module de calcul>} \\
Paramètre principal & $n = 100$ observations \\
\end{tabular}
\end{center}

\clearpage
\titrebandeau{Sommaire}
\markboth{Sommaire}{}
% \tableofcontents composerait son propre titre, en doublon avec le bandeau :
% on n'appelle que \@starttoc, qui compose les entrées seules.
\makeatletter
\@starttoc{toc}
\makeatother
\newpage


\partie{Partie I}{Identité visuelle}

% =============================================================================
\section{Palette et typographie}
\label{sec:palette}
% =============================================================================

\subsection{Couleurs}

La charte repose sur deux couleurs d'identité et deux gris de tableau. Le bleu
porte la structure (filets, titres de sous-section, liens, étiquettes), le
vert porte l'accent (titres de section et de partie, numéro de page, sorties).

\begin{center}
\small
\rowcolors{2}{tablerow}{white}
\begin{tabular}{@{}p{3.0cm}p{2.6cm}p{1.4cm}p{7.0cm}@{}}
\rowcolor{tableheader}
\textcolor{financeblue}{\textbf{Nom}} &
\textcolor{financeblue}{\textbf{RVB}} &
\textcolor{financeblue}{\textbf{Aperçu}} &
\textcolor{financeblue}{\textbf{Usage}} \\
\code{financeblue} & 0, 70, 140 & \textcolor{financeblue}{\rule{1cm}{2.2ex}} &
Filets, titres de sous-section, sommaire, liens, rubriques de fiche \\
\code{maincolor} & 0, 70, 140 & \textcolor{maincolor}{\rule{1cm}{2.2ex}} &
En-tête, pied de page, étiquettes de légende, références \\
\code{greenline} & 0, 180, 80 & \textcolor{greenline}{\rule{1cm}{2.2ex}} &
Titres de section et de partie, numéro de page, amorce de la page de garde \\
\code{tableheader} & 200, 220, 250 & \textcolor{tableheader}{\rule{1cm}{2.2ex}} &
Ligne d'en-tête des tableaux \\
\code{tablerow} & 245, 245, 245 & \textcolor{tablerow}{\rule{1cm}{2.2ex}} &
Lignes alternées des tableaux, fond des encadrés \\
\end{tabular}
\end{center}

\subsection{Typographie}

Le texte courant est composé en Termes (clone de Times), corps 11~pt,
interligne multiplié par 1,08\fsp; le code, en \code{txtt}, via la macro
\code{\textbackslash code\{\}}\fsp; les mathématiques restent en Computer Modern,
par exemple $\bar x = \frac{1}{n}\sum_{i=1}^{n} x_i$.

\subsubsection{Niveau de sous-sous-section}

Troisième niveau de titre, en gras bleu au corps du texte. Il n'apparaît pas
au sommaire (profondeur limitée à deux niveaux).

\subsection{Niveaux de titre}

\begin{itemize}[itemsep=2pt]
\item \textbf{Partie}~: \code{\textbackslash partie\{Partie I\}\{Titre\}},
  nouvelle page, filet bleu, mention en petites capitales, titre vert.
\item \textbf{Section}~: \code{\textbackslash section\{Titre\}}, bandeau
  (filet bleu épais, titre vert en 19~pt), numéro visible seulement au
  sommaire et dans l'en-tête.
\item \textbf{Sous-section et sous-sous-section}~: gras bleu, numérotées.
\item \textbf{Bandeau seul}~: \code{\textbackslash titrebandeau\{Titre\}}
  suivi de \code{\textbackslash markboth\{Titre\}\{\}} pour une page non
  numérotée (objet, sommaire).
\end{itemize}


\partie{Partie II}{Blocs de texte et tableaux}

% =============================================================================
\section{Encadrés, citations et listes}
\label{sec:blocs}
% =============================================================================

\subsection{Encadré générique}

L'encadré sert aux remarques de portée et aux extraits de code. Il se coupe
entre deux pages si nécessaire.

\begin{encadre}
\small
\code{donnees <- lire\_donnees("chemin/des/donnees")}\\
\code{res <- calculer(donnees, methode = "standard", n = 100,}\\
\code{\phantom{res <- calculer(}graine = 20260831)}\\
\code{res\$resultat\_final \ \ ;\ \ tableau\_resultats(res)}
\end{encadre}

\begin{encadre}
\textbf{Portée.} L'encadré de portée dit ce que l'outil fait et ce qu'il ne
fait pas~: les cas du texte de référence qu'il ne couvre pas, et les méthodes
qu'il met en \oe uvre.
\end{encadre}

\subsection{Citation réglementaire et référence}

\refl{A.~Auteur, \emph{Titre de l'article}, \emph{Nom de la revue}, vol.(numéro),
année, p.~1--26.}

\reglement{Annexe, section A, paragraphe 2}{Second extrait du texte de
référence.}

\subsection{Listes}

Liste à puces (tiret cadratin), avec un écart réglable par
\code{[itemsep=...]}~:

\begin{itemize}[itemsep=3pt]
\item premier élément, sur une ligne\fsp;
\item second élément, plus long, pour montrer l'alignement des lignes suivantes
  sous le texte de l'élément et non sous la puce\fsp;
  \begin{itemize}
  \item élément de second niveau.
  \end{itemize}
\item dernier élément.
\end{itemize}

Liste numérotée~:

\begin{enumerate}[itemsep=2pt]
\item contrôle des données d'entrée\fsp;
\item calcul\fsp;
\item restitution des résultats.
\end{enumerate}

% =============================================================================
\section{Tableaux}
\label{sec:tableaux}
% =============================================================================

\subsection{Tableau courant}

En-tête sur fond \code{tableheader} en gras bleu, lignes alternées
\code{tablerow} et blanc, aucun filet vertical.

\begin{center}
\small
\rowcolors{2}{tablerow}{white}
\begin{tabular}{@{}p{4.3cm}p{5.1cm}p{5.4cm}@{}}
\rowcolor{tableheader}
\textcolor{financeblue}{\textbf{Fichier}} &
\textcolor{financeblue}{\textbf{Contenu}} &
\textcolor{financeblue}{\textbf{Interdit}} \\
\code{<module de calcul>} & Logique métier, contrôles des données, calculs &
Toute primitive d'interface \\
\code{<point d'entrée>} & Saisie, déclenchement, agencement de l'interface &
Tout calcul métier \\
\code{<module d'affichage>} & Formatage des nombres, tableaux et tracés &
Tout calcul \\
\end{tabular}
\end{center}

\subsection{Tableau long en paysage}

Les tableaux longs utilisent \code{longtable} et les filets \code{booktabs},
en page paysage. Le tableau de la page suivante en montre la structure~:
en-tête répété, mention de suite, intertitres de groupe en gras bleu.

\begin{landscape}
\begin{center}
\scriptsize
\setlength{\tabcolsep}{4.5pt}
\renewcommand{\arraystretch}{1.22}
\begin{longtable}{@{}p{4.4cm}p{0.9cm}p{5.7cm}p{7.8cm}p{2.1cm}@{}}
\toprule
\textbf{Élément (fiche)} & \textbf{Réf.} & \textbf{Fonction(s) --- calcul} &
\textbf{Fonction(s) / méthode --- restitution} & \textbf{Clé} \\
\midrule
\endfirsthead
\toprule
\textbf{Élément (fiche)} & \textbf{Réf.} & \textbf{Fonction(s) --- calcul} &
\textbf{Fonction(s) / méthode --- restitution} & \textbf{Clé} \\
\midrule
\endhead
\multicolumn{5}{@{}l@{}}{\textit{Suite page suivante}} \\
\endfoot
\bottomrule
\endlastfoot

\multicolumn{5}{@{}l@{}}{\textbf{\textcolor{financeblue}{Groupe A --- intitulé du premier groupe}}} \\
\midrule
Premier élément & \ref{fiche:1} & \code{calcul\_a()} &
\code{restitution\_a()} (formule exacte) & \code{A1} \\
Deuxième élément & \ref{fiche:1} & \code{calcul\_b()} &
\code{restitution\_b()} (approximation, loi $\chi^2$) & --- \\
\midrule
\multicolumn{5}{@{}l@{}}{\textbf{\textcolor{financeblue}{Groupe B --- intitulé du second groupe}}} \\
\midrule
Troisième élément & \ref{fiche:1} & \code{calcul\_c()} &
\code{restitution\_c()} (simulation) & \code{C1} \\
\end{longtable}
\end{center}
\end{landscape}


\partie{Partie III}{Fiches, figures et notes de mise en \oe uvre}

% =============================================================================
\section{Fiche de test}
\label{sec:fiche}
% =============================================================================

Chaque test fait l'objet d'une fiche en six rubriques fixes, ouverte par
l'environnement \code{fiche} et la macro \code{\textbackslash Cadre}. Les
rubriques suivantes sont appelées par \code{\textbackslash Hzero},
\code{\textbackslash Stat}, \code{\textbackslash Loi}, \code{\textbackslash
Pval} et \code{\textbackslash Usage}.

\begin{fiche}{Intitulé de la fiche}
\refl{A.~Auteur, \emph{Titre de l'ouvrage}, éditeur, année, p.~10--12.}
\Cadre
Données, notations et hypothèses sous lesquelles le test est valide, sans être
vérifiées par lui~: un échantillon $x_1,\dots,x_n$ indépendant et de même loi.
\Hzero
$H_0 : \mu = \mu_0$.
\Stat
\[
  T = \frac{\bar x - \mu_0}{s/\sqrt{n}}.
\]
\Loi
Loi de $T$ sous $H_0$, exacte ou approchée, avec ses conditions.
\Pval
$p = 2\,\bigl(1 - F(|T_{\text{obs}}|)\bigr)$, $F$ fonction de répartition de la
loi précédente.
\Usage
Domaine d'emploi, sensibilité aux hypothèses et limites connues\fsp; renvoi à
la fonction du module de calcul qui le met en \oe uvre.
\end{fiche}

% =============================================================================
\section{Figure légendée}
\label{sec:figure}
% =============================================================================

Les figures sont placées à l'endroit exact de l'appel par l'environnement
\code{figureH} (et \code{tableH} pour un tableau légendé)\fsp; les environnements
flottants \code{figure} et \code{table} restent disponibles. L'étiquette de
légende est en gras, en petites capitales et en couleur \code{maincolor}\fsp;
une légende d'une ligne est centrée, une légende plus longue est justifiée
(figure~\ref{fig:exemple}).

\begin{figureH}
\centering
\fcolorbox{financeblue}{tablerow}{\begin{minipage}[c][4cm][c]{9cm}
\centering\small\color{financeblue} Emplacement de la figure\\
\code{\textbackslash includegraphics[width=...]\{fichier.pdf\}}
\end{minipage}}
\caption{Exemple de légende longue. Le texte de la légende est composé en
romain au corps du texte, justifié sur toute la largeur\fsp; seule l'étiquette
porte la couleur de la charte.}
\label{fig:exemple}
\end{figureH}

\begin{figureH}
\centering
\fcolorbox{financeblue}{tablerow}{\begin{minipage}[c][2cm][c]{9cm}
\centering\small\color{financeblue} Emplacement de la figure
\end{minipage}}
\caption{Légende courte, centrée.}
\end{figureH}

% =============================================================================
\section{Notes de mise en \oe uvre}
\label{sec:notes}
% =============================================================================

\subsection{Packages remplacés}

\begin{center}
\small
\rowcolors{2}{tablerow}{white}
\begin{tabular}{@{}p{3.0cm}p{11.8cm}@{}}
\rowcolor{tableheader}
\textcolor{financeblue}{\textbf{Package retiré}} &
\textcolor{financeblue}{\textbf{Remplacement}} \\
\code{newtxtext} & \code{tgtermes} (repli sur \code{ptm}) pour le texte,
\code{txtt} pour le code \\
\code{babel-french} & Césure française, \code{\textbackslash frenchspacing},
libellés, date, retrait du premier paragraphe, puces et espacement des
listes, réglés dans le préambule \\
\code{csquotes} & Guillemets « » saisis directement (codage T1) \\
\code{caption} & Redéfinition de \code{\textbackslash @makecaption} \\
\code{float} & Environnements \code{figureH} et \code{tableH} (bloc non flottant) \\
\code{enumitem} & \code{itemize} et \code{enumerate} redéfinis, clés
\code{itemsep}, \code{topsep}, \code{parsep}, \code{leftmargin}\fsp; liste
\code{rubriques} pour les fiches \\
\code{mdframed} & \code{framed}, avec cadre sur mesure \\
\code{titlesec} & \code{\textbackslash @startsection} et
\code{\textbackslash @seccntformat} \\
\code{tocloft} & Redéfinition de \code{\textbackslash l@section} \\
\code{pdflscape} & \code{lscape} et \code{/Rotate 90} ajouté aux pages paysage (\code{atbegshi}) \\
\code{tikz} & Filets de la page de garde superposés à la page par
\code{atbegshi} (sans crochet du noyau récent) \\
\code{dsfont}, \code{mathrsfs} & Non utilisés\fsp; \code{\textbackslash mathbb}
(amssymb) reste disponible \\
\end{tabular}
\end{center}

\subsection{Ponctuation}

Sans babel-french, l'espace qui précède les signes doubles n'est plus
ajoutée automatiquement, et une espace ordinaire saisie avant eux autorise
une coupure de ligne malvenue. On saisit donc \code{\textasciitilde:} pour le
deux-points (espace insécable) et \code{\textbackslash fsp;},
\code{\textbackslash fsp!}, \code{\textbackslash fsp?} pour les autres
signes (espace fine insécable)\fsp; les guillemets « » se saisissent
directement, avec des espaces ordinaires à l'intérieur.

\end{document}
